\def\SECTION{3.5 Equations of Parallel and Perpendicular Lines}
\documentclass[aspectratio=169,14pt]{beamer}
\usepackage[T1]{fontenc}
\usepackage{mathpazo}
\usepackage{amsmath,amssymb}
\usepackage{tikz}
\usetikzlibrary{arrows.meta,calc,decorations.pathreplacing}
\definecolor{fblue}{HTML}{1F4E79}
\definecolor{fteal}{HTML}{0F7C7C}
\definecolor{fgold}{HTML}{C9962B}
\definecolor{fred}{HTML}{B3261E}
\definecolor{fink}{HTML}{16233B}
\definecolor{fsoft}{HTML}{EEF3F8}
\definecolor{fgoldsoft}{HTML}{FBF1DC}
\setbeamertemplate{navigation symbols}{}
\setbeamercolor{normal text}{fg=fink}
\setbeamercolor{frametitle}{fg=fblue}
\setbeamerfont{frametitle}{series=\bfseries,size=\Large}
\setbeamercolor{structure}{fg=fteal}
\setbeamertemplate{itemize item}{\color{fteal}$\blacktriangleright$}
\setbeamertemplate{frametitle}{\vspace{6pt}\insertframetitle\par\vspace{-2pt}{\color{fgold}\rule{2.2cm}{2pt}}}
\setbeamertemplate{footline}{\hspace{0.6cm}{\color{fblue}\scriptsize\bfseries Fang\textcolor{fteal}{Math}}\hfill{\scriptsize\color{gray}\SECTION\quad\insertframenumber}\hspace{0.6cm}\vspace{6pt}}
\newcommand{\key}[1]{{\color{fteal}\textbf{#1}}}
\newcommand{\ans}[1]{\colorbox{fgoldsoft}{\color{fink}\ensuremath{#1}}}
\newcommand{\warn}[1]{\begin{beamercolorbox}[sep=8pt,rounded=true]{warnbox}{\color{fred}\textbf{Watch out!}}\enskip #1\end{beamercolorbox}}
\setbeamercolor{warnbox}{bg=fred!8,fg=fink}
\setbeamercolor{turnbox}{bg=fgoldsoft,fg=fink}
\setbeamercolor{defbox}{bg=fsoft,fg=fink}
\newcommand{\defn}[1]{\begin{beamercolorbox}[sep=8pt,rounded=true]{defbox}#1\end{beamercolorbox}}
\newenvironment{yourturn}{\begin{beamercolorbox}[sep=10pt,rounded=true]{turnbox}{\color{fgold}\Large\bfseries Your turn!}\quad{\small Pause the video and try it.}\par\medskip}{\end{beamercolorbox}}
\newcommand{\titleframe}[3]{%
\begin{frame}[plain]
\begin{tikzpicture}[remember picture,overlay]
\fill[fblue] (current page.north west) rectangle ([yshift=-0.9cm]current page.north east);
\node[anchor=west,text=white,font=\bfseries] at ([xshift=0.6cm,yshift=-0.45cm]current page.north west) {Fang\textcolor{fgold}{Math}};
\node[anchor=east,text=white,font=\small] at ([xshift=-0.6cm,yshift=-0.45cm]current page.north east) {Geometry $\cdot$ Chapter \CHNO};
\end{tikzpicture}
\vfill
{\color{fteal}\large\bfseries #1}\par\medskip
{\color{fblue}\fontsize{30}{34}\selectfont\bfseries #2}\par\medskip
{\color{fgold}\rule{3cm}{3pt}}\par\bigskip
{\large #3}
\vfill
{\small\color{gray} fangmath.com $\cdot$ YouTube @fangmathcom}
\end{frame}}
\newcommand{\closeframe}[1]{%
\begin{frame}{Remember}
#1
\vfill
\begin{center}{\color{fblue}\large\bfseries Now practice at your level:}\par\smallskip
{\Large\color{fteal}\bfseries fangmath.com} \enskip{\color{gray}$\to$ Geometry $\to$ Chapter \CHNO}\end{center}
\end{frame}}

\usepackage{pgfplots}\pgfplotsset{compat=1.17}
\newcommand{\grid}[4]{\draw[step=1,gray!25] (#1,#2) grid (#3,#4);\draw[-{Stealth},gray] (#1,0)--(#3+0.3,0) node[right]{\small$x$};\draw[-{Stealth},gray] (0,#2)--(0,#4+0.3) node[above]{\small$y$};}
\def\CHNO{3}
\usetikzlibrary{angles,quotes}
\begin{document}
\titleframe{Lesson 3.5}{Equations of Parallel \& Perpendicular Lines}{Same slope, or flip and change the sign}
\begin{frame}{Slope rules}
\onslide<1->{\key{Parallel} lines have the \key{same} slope.}\par\medskip
\onslide<2->{\key{Perpendicular} lines: slopes are \key{negative reciprocals}.}\par\medskip
\onslide<3->{Flip the fraction \textbf{and} change the sign: \ $\tfrac23 \to -\tfrac32$ \quad ($\tfrac23 \cdot -\tfrac32 = -1$)}\par\medskip
\end{frame}
\begin{frame}{Find the slopes}
\onslide<1->{$y = \tfrac23 x + 1$: \ slope $\tfrac23$}\par\medskip
\onslide<2->{Parallel slope: \ans{\tfrac23}}\par\medskip
\onslide<3->{Perpendicular slope: \ans{-\tfrac32}}\par\medskip
\end{frame}
\begin{frame}{Write a parallel line}
\onslide<1->{Through $(2, 3)$, parallel to $y = 4x - 1$}\par\medskip
\onslide<2->{Same slope $m = 4$: \ $y - 3 = 4(x - 2)$}\par\medskip
\onslide<3->{$y = \ans{4x - 5}$}\par\medskip
\end{frame}
\begin{frame}{Write a perpendicular line}
\onslide<1->{Through $(4, 1)$, perpendicular to $2x + 3y = 6$}\par\medskip
\onslide<2->{Solve for $y$: $y = -\tfrac23 x + 2$, \ slope $-\tfrac23$}\par\medskip
\onslide<3->{Perpendicular slope $\tfrac32$: \ $y - 1 = \tfrac32(x - 4)$ \ $\Rightarrow$ \ $y = \ans{\tfrac32 x - 5}$}\par\medskip
\end{frame}
\begin{frame}{Watch out}
\onslide<1->{\warn{Perpendicular needs \textbf{both} steps: flip \textbf{and} change the sign. The perpendicular slope of $2$ is $-\tfrac12$, not $\tfrac12$.}}\par\medskip
\end{frame}
\begin{frame}{}
\begin{yourturn}Write the line through $(3, -2)$ perpendicular to $y = -3x + 7$.\end{yourturn}\par\medskip
\onslide<2->{Slope $\tfrac13$: \ $y + 2 = \tfrac13(x - 3)$ \ $\Rightarrow$ \ $y = \ans{\tfrac13 x - 3}$}\par\medskip
\end{frame}
\closeframe{\begin{itemize}\setlength\itemsep{8pt}
\item \key{Parallel}: same slope.
\item \key{Perpendicular}: negative reciprocal --- flip and change the sign.
\item Then use point-slope form: $y - y_1 = m(x - x_1)$.
\end{itemize}}
\end{document}
